\section*{Chapitre \ref{ch:b1:crystals-ionic-covalent} --- Cristaux II~: solides ioniques, covalents et moléculaires}

\begin{solution}{exo:b1:crystals-ionic-covalent:1}
Anions~: $8 \times \frac18 + 6 \times \frac12 = 4$~; cations~: $12 \times
\frac14 + 1 = 4$. Chaque ion a 6 voisins de l'autre sorte (6:6). Quatre
cations et quatre anions~: \ce{NaCl}.
\end{solution}

\begin{solution}{exo:b1:crystals-ionic-covalent:2}
La moitié de la grande diagonale~: $a\sqrt3/2 = 412.3 \times 0.866 = \qty{357.1}{pm}$.
Un \ce{Cs+} et $8 \times \frac18 = 1$ \ce{Cl-}~: une unité formulaire.
\end{solution}

\begin{solution}{exo:b1:crystals-ionic-covalent:3}
$x = 102/181 = 0.56$, entre 0.414 et 0.732~: \omterm{def:b1:crystals-metals:coordination}{coordinence} octaédrique
(6), le \omterm{def:b1:crystals-ionic-covalent:structure-type}{type sel gemme}.
\end{solution}

\begin{solution}{exo:b1:crystals-ionic-covalent:4}
Métallique~: le cuivre. Ioniques~: le chlorure de sodium, le fluorure de
calcium. Covalents~: le diamant, le quartz. Moléculaires~: le diiode, la
glace.
\end{solution}

\begin{solution}{exo:b1:crystals-ionic-covalent:5}
Le long d'une arête, anion--cation--anion au contact~: $a = 2(r_+ + r_-)$.
Sur une diagonale de face se trouvent deux anions à la distance
$a\sqrt2/2$~; ils ne doivent pas s'interpénétrer~: $a\sqrt2/2 \ge 2r_-$,
donc $(r_+ + r_-)\sqrt2 \ge 2r_-$, c'est-à-dire
$r_+/r_- \ge \sqrt2 - 1$.
\end{solution}

\begin{solution}{exo:b1:crystals-ionic-covalent:6}
$\rho = 4 \times 58.44/(6.022 \times 10^{23} \times (5.641 \times
10^{-8})^3) = \qty{2.16}{g/cm^3}$, contre \qty{2.17}{g/cm^3} mesurés.
\end{solution}

\begin{solution}{exo:b1:crystals-ionic-covalent:7}
$\rho = 168.36/(6.022 \times 10^{23} \times (4.123 \times 10^{-8})^3) =
\qty{3.99}{g/cm^3}$, la valeur mesurée.
\end{solution}

\begin{solution}{exo:b1:crystals-ionic-covalent:8}
\ce{Zn-S} $= a\sqrt3/4 = \qty{234.2}{pm}$;
$\rho = 4 \times 97.44/(6.022 \times 10^{23} \times (5.409 \times
10^{-8})^3) = \qty{4.09}{g/cm^3}$, proche des \qty{4.04}{g/cm^3}
mesurés sur la sphalérite naturelle.
\end{solution}

\begin{solution}{exo:b1:crystals-ionic-covalent:9}
\ce{C-C} $= 245.6/\sqrt3 = \qty{141.8}{pm}$~; entre plans $669.6/2 =
\qty{334.8}{pm}$. Volume de la \omterm{def:b1:crystals-metals:lattice}{maille} $\frac{\sqrt3}{2}a^2c = \qty{3.498e-23}{cm^3}$,
$\rho = 4 \times 12.01/(6.022 \times 10^{23} \times 3.498 \times 10^{-23})
= \qty{2.28}{g/cm^3}$, contre \qty{3.52}{g/cm^3} pour le diamant~: les
plans sont très écartés, retenus seulement par des forces de London.
\end{solution}

\begin{solution}{exo:b1:crystals-ionic-covalent:10}
Voir la démonstration de la \cref{prop:b1:crystals-ionic-covalent:diamond}~:
$C = \pi\sqrt3/16 = 0.34$. La dureté vient de la force et de la rigidité
tridimensionnelle des \omterm{def:b1:lewis-resonance-vsepr:lewis}{liaisons covalentes}, non de l'empilement~: chaque
atome n'a que quatre voisins, à des angles fixés, ce qui laisse beaucoup
de vide.
\end{solution}

\begin{solution}{exo:b1:crystals-ionic-covalent:11}
\ce{Si-Si} $= a\sqrt3/4 = \qty{235.2}{pm}$, \omterm{def:b1:periodicity:radius}{rayon covalent}
\qty{117.6}{pm}~; $\rho = 8 \times 28.09/(6.022 \times 10^{23} \times
(5.431 \times 10^{-8})^3) = \qty{2.33}{g/cm^3}$, la valeur mesurée.
\end{solution}

\begin{solution}{exo:b1:crystals-ionic-covalent:12}
$x = 152/181 = 0.84 > 0.732$~: la règle prévoit le \omterm{def:b1:crystals-ionic-covalent:structure-type}{type chlorure de
césium}. Dans la structure sel gemme, \ce{Rb-Cl} $= a/2 = \qty{329.1}{pm}$,
proche de $152 + 181 = \qty{333}{pm}$~: la structure observée est
cohérente. La règle traite les ions comme des sphères dures et ignore
leur \omterm{def:b1:periodicity:polarisability}{polarisabilité} ainsi que la faible différence d'énergie entre les
deux structures (le chlorure de rubidium adopte le \omterm{def:b1:crystals-ionic-covalent:structure-type}{type chlorure de
césium} sous pression)~; c'est un guide, non une loi.
\end{solution}

\begin{solution}{pb:b1:crystals-ionic-covalent:1}
\textbf{1.} Aux 8 sommets et aux 6 centres des faces~: $8 \times \frac18 + 6
\times \frac12 = 4$ ions.
\textbf{2.} 8 \omterm{def:b1:crystals-metals:site}{sites tétraédriques}, tous occupés~: 8 \ce{F-} pour 4
\ce{Ca^{2+}}, rapport 1:2, \ce{CaF2}.
\textbf{3.} \ce{F-}~: 4 \ce{Ca^{2+}} aux sommets d'un tétraèdre.
\ce{Ca^{2+}}~: 8 \ce{F-} aux sommets d'un cube.
\textbf{4.} $a\sqrt3/4 = 546.3 \times 0.4330 = \qty{236.6}{pm}$.
\textbf{5.} $112 + 131 = \qty{243}{pm}$~: à 3\,\% près.
\textbf{6.} Les ions \ce{F-} forment un \omterm{def:b1:crystals-metals:lattice}{réseau} cubique simple d'arête
$a/2 =
\qty{273.2}{pm}$, un peu plus que $2 \times 131 = \qty{262}{pm}$~: ils se
touchent presque.
\textbf{7.} $x = 112/131 = 0.85$.
\textbf{8.} $x > 0.732$~: \omterm{def:b1:crystals-metals:coordination}{coordinence} 8.
\textbf{9.} Le \omterm{def:b1:crystals-ionic-covalent:structure-type}{type chlorure de césium} a autant de cations que d'anions~;
avec deux fois plus d'anions, seule la moitié des cubes d'anions peut
contenir un cation.
\textbf{10.} Les 8 \ce{F-} de la \omterm{def:b1:crystals-metals:lattice}{maille} forment un cube de 8 petits cubes
d'arête $a/2$~; les ions \ce{Ca^{2+}} occupent les centres de 4 d'entre
eux, en alternance, comme les cases noires d'un échiquier à trois
dimensions.
\textbf{11.} $C = \frac43\pi(4 \times 112^3 + 8 \times 131^3)/546.3^3 = 0.61$.
\textbf{12.} \ce{O^{2-}} sur le \omterm{def:b1:crystals-metals:lattice}{réseau} CFC (sommets et centres des
faces), \ce{Li+} dans les 8 \omterm{def:b1:crystals-metals:site}{sites tétraédriques}.
\textbf{13.} $a\sqrt3/4 = \qty{200.0}{pm}$~; $59 + 142 = \qty{201}{pm}$.
\textbf{14.} $\rho = 4 \times 29.88/(6.022 \times 10^{23} \times (4.619
\times 10^{-8})^3) = \qty{2.01}{g/cm^3}$, la valeur mesurée.
\textbf{15.} La blende~: le diamant est une blende dont les deux sortes
de positions sont occupées par du carbone.
\textbf{16.} $a\sqrt3/4 = \qty{154.5}{pm}$~; $\rho = 8 \times
12.01/(6.022 \times 10^{23} \times (3.567 \times 10^{-8})^3) =
\qty{3.52}{g/cm^3}$.
\textbf{17.} Dans le diamant, seule la moitié des \omterm{def:b1:crystals-metals:site}{sites tétraédriques} est
occupée, et par des atomes aussi gros que ceux du \omterm{def:b1:crystals-metals:lattice}{réseau}, qui ne se
touchent que le long des liaisons~: $C = 0.34$. Dans la fluorine, tous
les sites sont occupés et les petits cations comblent les vides entre de
gros anions~: $C = 0.61$.
\textbf{18.} $40.08 + 2 \times 19.00 = \qty{78.08}{g/mol}$.
\textbf{19.} $4 \times 78.08/(6.022 \times 10^{23}) = \qty{5.186e-22}{g}$.
\textbf{20.} $(5.463 \times 10^{-8})^3 = \qty{1.630e-22}{cm^3}$.
\textbf{21.} Un \ce{F-} manquant toutes les cent \omterm{def:b1:crystals-metals:lattice}{mailles} retire
$19.00/(100
\times 312.3)$ de la masse, soit 0.06\,\%~: invisible à trois chiffres
(et un \ce{Ca^{2+}} doit partir aussi, sinon le \omterm{def:b1:crystals-metals:crystal}{cristal} ne serait pas
neutre).
\textbf{22.} $\rho = 5.186 \times 10^{-22}/1.630 \times 10^{-22} =
\textbf{\qty{3.18}{g/cm^3}}$, exactement la masse volumique mesurée~: le
modèle de la \omterm{def:b1:crystals-metals:lattice}{maille}, construit à partir d'une seule mesure de $a$ aux
rayons X, rend compte de la masse du minéral.
\end{solution}
